\section{Machine Learning Surrogate Model Development}
\label{sec:ml}

Two surrogate modeling tasks were pursued: a scalar surrogate predicting the area-averaged Nusselt number, trained and evaluated on the full 5,460-case dataset described above (all cases converged; none were dropped, so no filtering step was required), and a spatial surrogate predicting the full temperature and velocity field at arbitrary points within the domain.

\subsection{Scalar Nusselt Number Surrogate}
\label{subsec:scalar_arch}

Three candidate architectures were trained and compared on equal footing, rather than an architecture being chosen in advance. A gradient-boosted tree ensemble (XGBoost) was trained with hyperparameters selected by five-fold cross-validated grid search over the number of estimators (200, 400, 800), maximum tree depth (3, 5, 7), and learning rate (0.03, 0.1); the selected configuration was 800 estimators, depth 5, learning rate 0.1 (cross-validation mean $R^2=0.9993$). A Kolmogorov-Arnold Network (KAN, width $[4,8,1]$, grid size 12, cubic splines) was trained with the LBFGS optimizer. A physics-informed neural network (PINN, a three-hidden-layer multilayer perceptron, 64 units each, $\tanh$ activation) was trained with the Adam optimizer with a soft penalty enforcing the two physically expected monotonic trends established in Section~\ref{subsec:governing}: $Nu$ non-decreasing in $Re$ at fixed $Ri$, and non-decreasing in $Ri$ at fixed $Re$ (both checked directly against the dataset before being encoded as constraints, not assumed). All three were trained and evaluated on an identical 70/30 train-test split (fixed random seed) of the four inputs $(R^\ast, \varepsilon, Re, Ri)$ against the single output $Nu$; $Re$ and $Ri$ were $\log_{10}$-transformed before normalization for the KAN and PINN candidates, reflecting the more-than-two-orders-of-magnitude range each spans.

\subsection{Spatial Field Surrogate}
\label{subsec:field_arch}

Three candidates were compared for the spatial field target, in the same architecture design as the initial study: a per-node multilayer perceptron conditioning each query point jointly on the case parameters and its own $(x,y)$ coordinate, $\mathbf{x} = [R^\ast,\ \varepsilon,\ Re,\ Ri,\ x,\ y]$, with output $[\theta(x,y),\ u(x,y),\ v(x,y)]$; a DeepONet encoding the case parameters through a branch network and the query coordinate through a trunk network, combined via a bilinear (dot-product) form; and a graph neural network (GNN) operating on the true finite-volume mesh connectivity extracted directly from the solver's mesh files (verified exact for all 35 geometries: every mesh node matched its field-export coordinate to within numerical precision, with zero unmatched nodes or dropped edges). Case-level, not node-level, splitting is used throughout (24 training geometries, 11 held-out test geometries), so that test performance measures generalization to entirely unseen geometries rather than interpolation within an already-seen mesh. The per-node field export (temperature and velocity at every mesh node, for all 5,460 cases, 13.1 million node-rows total) was generated by re-solving every case with the exact same physics and convergence criterion as the scalar dataset, imported directly from that dataset's own generation script rather than reimplemented; the resulting Nusselt number, recomputed independently during this second pass, matched the original scalar dataset exactly at all 5,460 points (maximum absolute difference $0.0$), the strongest available check that the two datasets describe the same underlying solutions.

\subsection{Training Protocol and Validation}
\label{subsec:training}

Beyond the holdout split, the scalar surrogate's winning architecture was further validated by Leave-One-Geometry-Out (LOGO) cross-validation across all 35 geometries, retraining 35 times with one distinct $(R^\ast,\varepsilon)$ combination withheld entirely each time, and, critically, against fresh numerical simulations run at $(Re,Ri)$ combinations chosen to fall strictly between grid levels on both axes, mathematically guaranteed to be absent from the 5,460-case sampling grid, providing a genuine measure of out-of-sample accuracy rather than in-grid interpolation accuracy.

Model performance was assessed using the coefficient of determination $R^2$, root mean square error (RMSE), and mean absolute percentage error (MAPE), computed in the original $Nu$ units.

\subsection{Uncertainty Quantification}
\label{subsec:uq}

For the winning scalar architecture, a tree ensemble rather than a neural network, uncertainty was quantified with a 20-member bootstrap ensemble (each member trained on an independent bootstrap resample of the training set with a different random seed), the natural analogue of a deep ensemble for a non-neural architecture. For the winning spatial-field architecture, a true deep ensemble of five independently initialized networks was used. Calibration was assessed as the fraction of true values falling within each ensemble's own reported 95\% predictive interval, evaluated both on the same-grid test split and on the true out-of-sample points described above.

\subsection{Interpretability}
\label{subsec:interp}

SHAP (SHapley Additive exPlanations) values were computed for the winning scalar surrogate on its held-out test set to rank each input's relative contribution to the predicted Nusselt number.
